\subsection{TENSOR PRODUCT OF MATRICES OVER A COMMUTATIVE RING}

Let C be a commutative ring, E, F, U, V four C-modules and \(\phi:E\to U\) , \(\psi:F\to V\) two C-linear mappings. Suppose that E, F, U, V have respectively finite bases \((e_{\lambda})_{\lambda \in L}, (f_{\mu})_{\mu \in M}, (u_{\rho})_{\rho \in R}, (v_{\sigma})_{\sigma \in S}\) ; let \(A = (a_{\rho \lambda})\) be the matrix of \(\phi\) with respect to \((e_{\lambda})\) and \((u_{\rho})\) , \(B = (b_{\sigma \mu})\) that of \(\psi\) with respect to \((f_{\mu})\) and \((v_{\sigma})\) . For every ordered pair \((\lambda, \mu) \in L \times M = N\) , let \(g_{\lambda \mu} = e_{\lambda} \otimes f_{\mu}\) ; for every ordered pair \((\rho, \sigma) \in R \times S = T\) let \(w_{\rho \sigma} = u_{\rho} \otimes v_{\sigma}\) ; the \(g_{\lambda \mu}\) then form a basis of \(E \otimes F\) and the \(w_{\rho \sigma}\) a basis of \(U \otimes V\) ( \(\S 3\) , no. 6, Corollary 2 to Proposition 7). The tensor product of \(A\) by \(B\) , denoted by \(A \otimes B\) , is the matrix \(X = (x_{\tau v})(\tau, v) \in T \times N\) whose elements are given by

\[
x _ {(\rho , \sigma), (\lambda , \mu)} = a _ {\rho \lambda} b _ {\sigma \mu}.\tag{42}
\]

Then \(A \otimes B\) is the matrix of \(\phi \otimes \psi\) with respect to the bases \((g_{\lambda \mu})\) and \((w_{\rho \sigma})\) . By definition (§ 3, no. 2, formula (3))

\[
\begin{array}{r l} (\phi \otimes \psi) (g _ {\lambda \mu}) & = (\phi \otimes \psi) (e _ {\lambda} \otimes f _ {\mu}) = \phi (e _ {\lambda}) \otimes \psi (f _ {\mu}) \\ & = \sum_ {\rho , \sigma} a _ {\rho \lambda} b _ {\sigma \mu} (u _ {\rho} \otimes v _ {\sigma}) = \sum_ {\rho , \sigma} a _ {\rho \lambda} b _ {\sigma \mu} w _ {\rho \sigma}. \end{array}
\]

Definition (42) of the elements of \(A \otimes B\) shows that this matrix corresponds bijectively with the matrix of matrices \((a_{\rho \lambda}B)_{(\rho, \lambda) \in \mathbb{R} \times \mathbb{L}}\) and also the matrix \((Ab_{\sigma \mu})_{(\sigma, \mu) \in \mathbb{S} \times \mathbb{M}}\) (no. 5).

The fact that \((\phi, \psi) \mapsto \phi \otimes \psi\) is a C-bilinear mapping and formula (9) of §3, no. 5, can be expressed by the identities

\[
\left\{ \begin{array}{l} A \otimes (B _ {1} + B _ {2}) = A \otimes B _ {1} + A \otimes B _ {2} \\ (A _ {1} + A _ {2}) \otimes B = A _ {1} \otimes B + A _ {2} \otimes B \end{array} \right.\tag{43}
\]

\[
(c A) \otimes B = A \otimes (c B) = c (A \otimes B) \quad \text { for } c \in \mathbf {C}\tag{44}
\]

\[
(A _ {1} \otimes B _ {1}) (A _ {2} \otimes B _ {2}) = (A _ {1} A _ {2}) \otimes (B _ {1} B _ {2})\tag{45}
\]

when the operations appearing are defined. The transpose of a tensor product of matrices is given by

\[
{ } ^ { t } ( A \otimes B ) = ( { } ^ { t } A ) \otimes ( { } ^ { t } B ) .\tag{46}
\]

If A and B are invertible square matrices over C, \(A \otimes B\) is invertible and

\[
(A \otimes B) ^ {- 1} = (A ^ {- 1}) \otimes (B ^ {- 1}).\tag{47}
\]

Let \((e_{\lambda}^{\prime})_{\lambda \in L}\) be another basis of E and \((f_{\mu}^{\prime})_{\mu \in M}\) another basis of F; if P is the matrix of passage from the basis \((e_{\lambda})\) to the basis \((e_{\lambda}^{\prime})\) and Q the matrix of passage from the basis \((f_{\mu})\) to the basis \((f_{\mu}^{\prime})\) , the matrix of passage from the basis \((e_{\lambda} \otimes f_{\mu})\) to the basis \((e_{\lambda}^{\prime} \otimes f_{\mu}^{\prime})\) is \(P \otimes Q\) . If \(A^{\prime}\) is equivalent (resp. similar) to A and \(B^{\prime}\) equivalent (resp. similar) to B, then \(A^{\prime} \otimes B^{\prime}\) is equivalent (resp. similar) to \(A \otimes B\) .

The definition of tensor product of matrices can be generalized in an obvious way to an arbitrary finite number of matrices over C; in particular we have the associativity formula

\[
\left(\bigotimes_ {i \in I _ {1}} X _ {i}\right) \otimes \left(\bigotimes_ {i \in I _ {2}} X _ {i}\right) = \bigotimes_ {i \in I} X _ {i}\tag{48}
\]

for every partition \((\mathbf{I}_{1}, \mathbf{I}_{2})\) of the finite indexing set I.

